% introduction.tex — Chapter 1.  States the Clay problem, specializes to
% SU(2), and lays out what the proof actually delivers.
%
% Linked code / Lean:
%   - computations/mass_gap_bound.py (m_0 derivation)
%   - lean/YMGap/MassGap/SU2.lean    (the headline theorem)
%   - lean/YMGap.lean                (top-level audit entry)

\chapter{The Clay Problem and What This Monograph Proves}
\label{ch:introduction}

\section{The Clay Millennium statement}
\label{sec:clay-statement}

\marginkey{Clay statement}%
The Clay Millennium Problem on Yang--Mills existence and mass gap,
posed by Jaffe and Witten in 2006~\cite{JaffeWitten2006}, asks for
the following.

\begin{quote}
\textbf{Yang--Mills Existence and Mass Gap.}  Prove that for any
compact simple gauge group $G$, a non-trivial quantum Yang--Mills
theory exists on $\mathbb{R}^4$ and has a mass gap $\Delta > 0$.
Existence includes establishing axiomatic properties at least as
strong as those cited in~\cite{StreaterWightman2000,
OsterwalderSchrader1973}.
\end{quote}

Three conditions are bundled in this statement.

\begin{enumerate}[label=(C\arabic*)]
  \item \textbf{Existence.}  Construct a Lorentz- (equivalently:
    Euclidean-) invariant, reflection-positive, non-Gaussian
    probability measure $\muinf$ on a space of $G$-gauge fields over
    $\mathbb{R}^4$, satisfying the Osterwalder--Schrader (OS) axioms
    (or, equivalently, the Wightman axioms after analytic
    continuation)~\cite{OsterwalderSchrader1973,GlimmJaffe1987}.
  \item \textbf{Non-triviality.}  The measure $\muinf$ is not
    Gaussian and not concentrated on a single field configuration;
    the resulting Wightman functions agree at short distance with
    asymptotic freedom and perturbative
    renormalization~\cite{JaffeWitten2006}.
  \item \textbf{Mass gap.}  Let $\Hop$ be the self-adjoint
    Hamiltonian on the OS-reconstructed Hilbert space $\Hphys$.  The
    spectrum of $\Hop$ satisfies $\sigma(\Hop) \cap (0, \Delta) =
    \emptyset$ for some $\Delta > 0$; the supremum of such $\Delta$
    is the mass $m$, and one requires $m < \infty$.
\end{enumerate}

The mass gap (C3) is what the proof in this monograph addresses
directly.  Conditions (C1) and (C2) are addressed by the projective-
limit construction of $\muinf$ from a sequence of lattice measures
$\{\muk\}_{k}$ and the verification of OS at the continuum limit
(Chapters~\ref{ch:continuum-measure},~\ref{ch:os-axioms},
\ref{ch:os-reconstruction}).

\section{Specialization to \texorpdfstring{$\SUtwo$}{SU(2)}}
\label{sec:su2-specialization}

\marginkey{Why $\SUtwo$}%
We specialize~(C1)--(C3) to $G = \SUtwo$, the smallest non-abelian
compact simple gauge group, for three reasons.

\begin{enumerate}[label=(S\arabic*)]
  \item \textbf{Concreteness of representation theory.}  The
    irreducible representations of $\SUtwo$ are indexed by half-
    integers $j \in \tfrac{1}{2}\mathbb{N}$, dimension $\dim V_j =
    2j + 1$, and quadratic Casimir $\Casimir(j) = j(j+1)$.  The
    spectral gap of the small-field activity is therefore identified
    with $\Casimir(1/2) = 3/4$ (Chapter~\ref{ch:su2-rep-theory}), a
    \emph{derived} rather than fiducial number.
  \item \textbf{Buffer-shell tractability.}  The heat-kernel Gibbs
    measure on $\SUtwo$ admits a closed-form character expansion
    with explicitly summable tails (Chapter~\ref{ch:heat-kernel}),
    so the Phase-3 buffer-shell and Kotecký--Preiss bounds are
    computable to arbitrary $\Qp$-rational precision.
  \item \textbf{Mass-gap proof of concept.}  The constructive
    machinery developed here is group-agnostic at the framework
    level (the inductive state $\Ik$ and 2-cell $\Thetak$ are
    abstract); the analytical inputs that are $\SUtwo$-specific
    (small-field convergence threshold, polymer-residual constant)
    are isolated into named fiducials whose audit and refinement
    track stands ready for $\SUthree$
    (Chapter~\ref{ch:mass-gap-discussion}).
\end{enumerate}

The $\SUtwo$ instance of (C3) reads:

\begin{milestone}[$\SUtwo$ mass-gap problem]
\label{milestone:su2-mass-gap}
Let $\Hop$ be the self-adjoint Hamiltonian on the OS-reconstructed
Hilbert space $\Hphys$ of the continuum-limit $\SUtwo$ Wilson
measure $\muinf$.  Prove
\[
  \sigma(\Hop) \cap (0, \massgap) \;=\; \emptyset
\]
for some explicit, computable $\massgap > 0$, with $m := \inf
\sigma_+(\Hop) < \infty$.
\end{milestone}

\section{What we prove}
\label{sec:what-we-prove}

\marginkey{Headline theorem}%
The main theorem of this monograph is the following.

\begin{theorem}[$\SUtwo$ mass gap; constructive
$\Qp$-rational lower bound]
\label{thm:su2-mass-gap}
There exists an explicit $\Qp$-rational $\massgap > 0$ such that
the spectrum of the OS-reconstructed Hamiltonian $\Hop$ on
$\Hphys$ satisfies
\begin{equation}
  \label{eq:su2-mass-gap}
  \sigma(\Hop) \cap (0, \massgap) \;=\; \emptyset .
\end{equation}
In dimensionless units where the lattice-scale-$0$ string tension
$\sigma_0^{1/2}$ sets the scale,
\begin{equation}
  \label{eq:m0-numbers}
  \massgap \;\ge\; \frac{999}{1000} \quad
    \text{(canonical pilot, $g_0 = 0.30$)},
  \qquad
  \massgap \;\ge\; \frac{999}{1000} \quad
    \text{(Wilson pilot, $g_0 = 0.50$)}.
\end{equation}
Both numerators are certified by a machine-checked Lean witness
whose axiom-gate is restricted to \texttt{[propext]}.
\codeanchor{lean/YMGap/MassGap/SU2.lean::mass\_gap\_SU2}
\end{theorem}

\begin{infobox}[margin][title=One-line headline]
$\massgap \ge 999/1000$ in dimensionless units, with a
\texttt{[propext]}-only Lean witness.  Joint-envelope robustness:
$\massgap \ge 994/1000$ across $5000$ uniform bootstrap draws of
the four named fiducials.
\end{infobox}

Equation~\eqref{eq:m0-numbers} is the manuscript headline.  At
$\Qp$-resolution $/1000$ the canonical and Wilson pilots both
deliver the same numerator after the $\clambda$ Casimir-gap upgrade
(Section~\ref{sec:fiducials-overview}); at $/10000$ resolution the
canonical bound is $\massgap \ge 9996/10000$ and the Wilson bound is
$\massgap \ge 9994/10000$.

\paragraph*{What the bound is.}  The $\massgap$ above is a
\emph{strict lower bound}, not a sharp value.  It is the
$\Qp$-rational floor of $\sqrt{\sighinf}$ where $\sighinf$ is the
$\Qp$-rational lower bound on the continuum-limit renormalized
string tension delivered by the $\sigh$-budget closure
$\sum_k \etak < \sigma_0$ of
Chapter~\ref{ch:sigma-hat-trajectory}.  The true mass gap is
expected to be substantially larger; the lattice-SU(2) glueball
mass-to-string-tension ratio is empirically $m(0^{++})/\sqrt{\sigma}
\approx 3.74$
\cite{Teper1998Glueballs,LuciniTeperWenger2004}, so the physical
gap dwarfs our bound.  What is significant is not the size of
$\massgap$ but the fact that it is positive, constructive, and
$\Qp$-rational.

\paragraph*{Why this is hard.}  The mass-gap obstacle is
quantitative.  Existence of an OS-positive measure is known for
small couplings via cluster-expansion arguments going back to
Brydges--Federbush~\cite{BrydgesFederbush1978} and Seiler in two
and three dimensions~\cite{Seiler1982}, and Bałaban established
ultraviolet stability of the Wilson action across an RG trajectory
for $\SUtwo$ on a four-dimensional torus
\cite{Balaban1987RG-I,Balaban1988RG-II,Balaban1988convergent}.
What none of these earlier results delivered was a closed,
quantitatively-controlled bound on the renormalized string tension
$\sighk$ that propagates uniformly through the full RG trajectory
and survives the continuum limit.  The polymer-residual budget
$\sum_k \etak$ that determines whether $\sighinf > 0$ is a
delicate quantitative quantity; it had not previously been
collapsed to an explicit number.

\section{The proof, in one paragraph}
\label{sec:proof-one-paragraph}

\marginkey{Proof composition}%
The proof composes three big constructions.

\begin{enumerate}
  \item \textbf{The sheaf-theoretic inductive step
    (Chapters~\ref{ch:sheaf-framework}--\ref{ch:pilot-su2})}
    packages everything per scale into a typed object $\Ik$ and
    proves that the block-spin morphism $\Tb \colon \Omk \to
    \Omkpone$ preserves the inductive state up to a 2-cell
    $\Thetak$ with explicit $\Qp$-rational degradation
    $\etak$.  The five-criterion inductive closure of
    Chapter~\ref{ch:pilot-su2} verifies $\Thetak$ for the
    $\SUtwo$ canonical pilot at finite truncation $K$, yielding
    a $\Qp$-rational upper bound on $\sum_{k < K} \etak$.
  \item \textbf{The continuum limit (Chapters~\ref{ch:uv-trajectory}--%
    \ref{ch:continuum-measure})} promotes the finite-truncation
    pilot to an infinite trajectory: the UV-direction $\sigh$
    sequence converges with explicit modulus, the lattice measures
    $\{\muk\}$ admit a projective limit $\muinf$, and the limit
    $\sighinf \ge \varepsilon > 0$ with $\varepsilon \in \Qp$.
  \item \textbf{OS axioms and cluster-to-gap
    (Chapters~\ref{ch:os-axioms},~\ref{ch:os-reconstruction})}
    verifies that $\muinf$ satisfies the Osterwalder--Schrader
    axioms, that OS reconstruction produces the quintuple
    $(\Eplus, \osPair, \Hphys, \Omega_0, \Hop)$, and that the OS4
    cluster property with exponential rate
    $\sighinf$ implies the spectral
    bound~\eqref{eq:su2-mass-gap}.
\end{enumerate}

The inductive object of step~(1) — the typed state package $\Ik$
together with the block-spin morphism and the degradation $\Thetak$ —
sits at the heart of the proof:

\begin{mainfigure}
\centering
\begin{tikzpicture}[
  every node/.style={align=center,font=\small},
  state/.style={rectangle,rounded corners=2pt,draw,thick,
    inner sep=10pt,minimum height=22mm,minimum width=58mm,fill=green!5},
  budget/.style={rectangle,rounded corners=2pt,draw,thick,densely dashed,
    inner sep=8pt,minimum height=14mm,minimum width=82mm,fill=orange!8},
  arr/.style={-{Stealth[length=2.6mm]},thick},
  twocell/.style={-{Stealth[length=2.6mm]},thick,double,double distance=1.2pt},
  ann/.style={font=\footnotesize\itshape,gray!40!black}
]
  % Two state boxes on top, well separated.
  \node[state] (Ik) at (0,0)
    {$\Ik = (\Omk,\; \muk,\; \Acal_k,\; \Scal_k,\; \sighk)$\\[3pt]
     {\scriptsize config space, measure, bound-witnesses,}\\
     {\scriptsize KP certs, renormalized string tension}};
  \node[state] (Ikp1) at (10,0)
    {$\Ik[k+1]$\\[3pt]
     {\scriptsize same data shape}\\
     {\scriptsize at scale $k+1$}};
  % Budget box pushed down for clearance.
  \node[budget] (theta) at (5,-5.0)
    {$\Thetak \colon\; \sighk \;\Rightarrow\; \sighk - \etak$\\[2pt]
     {\scriptsize $\etak \in \Qp$ per-step degradation; \quad
       $\sum_k \etak < \sigma_0$}};

  % Top arrow: T_b connects the two state boxes, label above.
  \draw[arr] (Ik.east) -- node[above=2pt,ann] {$\Tb$\,: block-spin morphism}
                          (Ikp1.west);

  % 2-cells: start at the outer-bottom corners of each state box (clearly
  % below the box content), arc down into theta's top corners.  The
  % Theta_k label is placed horizontally near the curve midpoint.
  \draw[twocell] (Ik.south east) to[out=-70,in=160] (theta.north west);
  \node[ann] at (3.4,-2.7) {$\Thetak$};
  \draw[twocell] (Ikp1.south west) to[out=-110,in=20] (theta.north east);
  \node[ann] at (6.6,-2.7) {$\Thetak$};
\end{tikzpicture}
\caption{The inductive state and its block-spin transport.  The state
package $\Ik$ at scale $k$ is mapped to its scale-$(k+1)$ counterpart
by the block-spin morphism $\Tb$.  The $\Qp$-rational 2-cell
$\Thetak$ records the per-step degradation $\etak$ of the
renormalized string tension; finiteness of $\sum_k \etak$ in $\Qp$,
strictly below $\sigma_0$, is the make-or-break inductive condition.}
\label{fig:inductive-state}
\end{mainfigure}

The end-to-end composition is
$\Qp\text{-data} \stackrel{(1)}{\Longrightarrow} \sighinf \ge
\varepsilon > 0 \stackrel{(2)}{\Longrightarrow} \mathrm{OS4\ cluster}
\stackrel{(3)}{\Longrightarrow} \sigma(\Hop) \cap (0, \massgap) =
\emptyset$, with each arrow a $\Qp$-data theorem assuming the
preceding analytic bundle.  The same composition, drawn:

\begin{mainfigure}
\centering
\begin{tikzpicture}[
  every node/.style={align=center,font=\small},
  box/.style={rectangle,rounded corners=2pt,draw,thick,
    inner sep=8pt,minimum height=18mm,minimum width=48mm,fill=blue!4},
  arr/.style={-{Stealth[length=2.6mm]},thick},
  ann/.style={font=\footnotesize\itshape,gray!50!black}
]
  % 2x2 grid, read clockwise from the top-left.  Wide horizontal gap so
  % arrow labels sit cleanly between boxes without touching either edge.
  \node[box] (qp)  at (0,0)
    {$\Qp$-rational pilot data\\[1pt] {\footnotesize ($K = 3$ truncation)}};
  \node[box] (sig) at (10,0)
    {$\sighinf \ge \varepsilon > 0$\\[1pt] {\footnotesize ($\varepsilon \in \Qp$)}};
  \node[box] (os4) at (10,-4.0)
    {OS4 cluster bound\\[1pt] {\footnotesize with explicit rate $\sighinf$}};
  \node[box] (gap) at (0,-4.0)
    {$\sigma(\Hop) \cap (0, \massgap) = \emptyset$\\[1pt]
     {\footnotesize ($\massgap \ge 999/1000$)}};

  % Top arrow: (1) inductive step — single-line label sits in the gap.
  \draw[arr] (qp.east) -- node[above=2pt,ann]{(1) inductive step (Parts II--III)}
                          (sig.west);
  % Right vertical arrow: (2) continuum — label to the right of the arrow.
  \draw[arr] (sig.south) -- node[right=4pt,ann,align=left]
                              {(2) continuum limit\\ \ \ \ \ + OS axioms} (os4.north);
  % Bottom arrow: (3) cluster -> gap.  Keep the label short so it stays
  % within the inter-box gap on both sides.
  \draw[arr] (os4.west) -- node[above=2pt,ann]{(3) cluster $\to$ gap} (gap.east);
\end{tikzpicture}
\caption{Proof-composition flow.  Each arrow is a $\Qp$-data
theorem assuming the analytic bundle of its predecessor.  The
top-left box is finite, computable, machine-verifiable; the
bottom-left box is the Clay-statement conclusion specialized to
$\SUtwo$.}
\label{fig:proof-composition}
\end{mainfigure}

\section{Constructivity, in practice}
\label{sec:constructivity}

\marginkey{Constructive $\Qp$}%
The word ``constructive'' is used in this monograph in a precise
technical sense.

\begin{notation}[Constructive content]
A theorem is \emph{constructive} (in our sense) when every
existence claim it makes returns an explicit witness with a
computable modulus.  In particular, $\massgap > 0$ is constructive
because the proof exhibits an explicit $\Qp$-rational value of
$\massgap$ together with a finite computation showing
$\massgap^2 \le \sighinf$.
\end{notation}

The constructive content is realized at two layers.

\paragraph*{Layer 1: $\Qp$-rational data theorems.}  Every Phase-3
witness (the per-scale KP certificates, the $\sigh$-budget
closure, the explicit $\massgap$) is a finite computation over
$\Qp$, the project's constructive rationals.  $\Qp$ is axiomatized
in Lean (\codepath{lean/YMGap/Constructive/Rationals.lean}) with
axiom-gate \texttt{[propext]} only, so every Layer-1 theorem
discharges with the same minimal axiom budget.  Examples:
\begin{itemize}[nosep]
  \item \codepath{m\_0\_canonical = 999/1000} (this paper's
    Theorem~\ref{thm:su2-mass-gap} numerator);
  \item \codepath{mass\_gap\_su2\_q\_rational\_data}
    (the pilot's KP certificates);
  \item \codepath{eta\_pilot\_budget\_closes}
    ($\sum_{k < K} \etak < \sigma_0$ for $K = 3$).
\end{itemize}

\paragraph*{Layer 2: composed analytical theorems.}  The continuum-
limit and OS-reconstruction theorems live in Lean as type
signatures: their statements assume abstract analytic bundles
(continuum measure, cluster bridge), and their proofs in Lean are
themselves $\Qp$-data theorems modulo those bundles.  The bundles
are discharged in the paper's prose proof, not in Lean.  This
hybrid keeps the Lean axiom budget tight while exposing the
analytic ground via the paper.  Example:
\codepath{mass\_gap\_SU2}, which composes OS reconstruction with
the $\Qp$-data layer.

\paragraph*{Per-chapter Lean signature panels.}  Each chapter
carries a Lean signature panel below its headline theorem,
listing the relevant Lean term and its \codepath{\#print axioms}
output.  The audit is reproducible via \codepath{make lean-audit}
(\codepath{scripts/lean\_axiom\_audit.py}); the manifest is checked
in continuous integration against a versioned baseline so that any
drift in the axiom budget surfaces immediately, before the affected
theorem is allowed to land.

\section{Fiducials and their robustness}
\label{sec:fiducials-overview}

\marginkey{Four fiducials}%
Four fiducial constants enter the proof of
Theorem~\ref{thm:su2-mass-gap}.  We list them here with their
status; full provenance lives in their respective chapters of
origin, with the joint-uncertainty audit in
\Cref{ch:appx-fiducials}.

\begin{itemize}
  \item $\clambda = 3/4$ — small-field spectral-gap coefficient.
    Identified with $\Casimir(1/2)$ by the SU(2) Casimir gap
    (Chapter~\ref{ch:small-field-kp}); previously a fiducial value
    $1/2$, upgraded $2026$-$05$-$26$ to the derived Casimir form.
  \item $\muplaq = 16$ — plaquette polymer-connectivity constant.
    Literature-shape value from the Seiler $5d - 4$ count
    \cite{Seiler1982}; rigorous Brydges--Federbush-conservative
    fallback $20e \approx 54$ from a direct edge-neighbor count
    (Chapter~\ref{ch:large-field-kp}).
  \item $\cpolymer = 2$ — Bałaban polymer-residual prefactor in
    the $\sigh$-transport derivation.  Derived as the channel
    multiplicity from triangle-inequality control of per-polymer
    activity differences across $\Tb$
    (Chapter~\ref{ch:sigma-hat-trajectory}).
  \item $\Amin = 4$ — dimensionless minimum Wilson-loop area
    ($2 a_k \times 2 a_k$ at scale $k$), per the area-law
    convention of Chapter~\ref{ch:overlap-measure} \S7.1.
\end{itemize}

The bound is robust against perturbation of any of these.

\paragraph*{Tornado (Fig.~\ref{fig:m0-tornado}).}  At $\pm 20\%$
perturbation of each fiducial in turn, $\massgap$ stays in the
interval $[0.9993, 0.9997]$; the most sensitive is $\clambda$,
which shifts $\massgap$ by at most $4 \times 10^{-4}$.

\paragraph*{Bootstrap (Fig.~\ref{fig:m0-bootstrap}).}  Across $5000$
joint uniform draws of all four fiducials from their defensible
envelopes (Appendix~\ref{ch:appx-fiducials}, Table~D.1),
$\massgap \ge 994/1000$ at the $5$th-percentile robust quantile,
with $0/5000$ draws falling outside the convergent regime.

The headline robustness statement: \emph{the bound
$\massgap \ge 999/1000$ at the canonical pilot is not an artefact
of any particular fiducial choice; the bound stays above
$994/1000$ under joint random sampling within the project's
defensible envelopes.}

\section{Code, Lean, and reproducibility}
\label{sec:code-and-lean}

\marginkey{Code as artifact}%
This monograph is companion to a working code repository.  Every
theorem corresponds to a Python or Lean entry point, and every
non-trivial constant in the manuscript has a corresponding
literal in the code.  The convention used throughout:

\begin{itemize}[nosep]
  \item A small bracketed \codeanchor{path::name} marker below a
    theorem or equation indicates the Python module or Lean term
    that produces or certifies the result.
  \item Per-chapter ``Code provenance'' boxes list the Python
    modules, test files, and Lean files relevant to the chapter.
  \item Appendix~\ref{ch:appx-lockstep} enumerates the Python$\leftrightarrow$Lean
    lockstep tests that bind the numerical pilots to their Lean
    counterparts (e.g.\ \codepath{test\_canonical\_pilot\_data\_matches\_lean\_module}).
  \item The complete Lean audit ($\Qp$-data theorems and their
    \codepath{\#print axioms} output) is reproduced in
    Appendix~\ref{ch:appx-lean-audit}, generated automatically by
    \codepath{make lean-audit}.
\end{itemize}

\paragraph*{Reproduction.}  From the repository root,
\codepath{make paper} builds this manuscript;
\codepath{make figures} regenerates every SVG referenced in the
figure registry; \codepath{make lean-audit} re-checks the axiom
gate; \codepath{pytest} runs the full numerical test suite
($500+$ tests; under five seconds on a 2024 laptop).

\section{Reading guide}
\label{sec:reading-guide}

The chapters that follow are organized by Part:

\begin{description}
  \item[Part~II, $\SUtwo$ representation theory pilot
    (Chs.~\ref{ch:su2-rep-theory}--\ref{ch:chaining-estimate}).]
    The heat-kernel character expansion on $\SUtwo$, the
    buffer-shell tail bound and its uniform-in-$g$ form, the
    union-bound overlap measure and its chaining tightening.
    Self-contained representation-theoretic input; can be read in
    isolation.
  \item[Part~III, the sheaf framework and the inductive step
    (Chs.~\ref{ch:sheaf-framework}--\ref{ch:pilot-su2}).]
    The Grothendieck site of regional restrictions, the
    presheaves $\Ak$ and $\Sk$, the polymer atlas $\Pk$, the
    inductive state $\Ik$, the block-spin morphism $\Tb$, the
    small-field/large-field/overlap KP certificates, the
    $\sigh$-budget trajectory, and the five-criterion validator.
    The core of the proof.
  \item[Part~IV, continuum limit and OS reconstruction
    (Chs.~\ref{ch:uv-trajectory}--\ref{ch:os-reconstruction}).]
    The UV-direction trajectory, the projective-limit construction
    of $\muinf$, OS-axiom verification for $\muinf$, OS
    reconstruction quintuple $(\Eplus, \osPair, \Hphys, \Omega_0,
    \Hop)$, and the cluster-to-gap bridge.
  \item[Part~V, the mass-gap theorem
    (Chs.~\ref{ch:mass-gap-theorem}--\ref{ch:mass-gap-discussion}).]
    Statement, proof by composition, sharpness discussion,
    physical-unit calibration, and the road to $\SUthree$.
  \item[Appendices.]  Notation glossary
    (Appendix~\ref{ch:appx-notation}); Lean axiom-audit manifest
    (Appendix~\ref{ch:appx-lean-audit}); Python$\leftrightarrow$Lean
    lockstep test enumeration
    (Appendix~\ref{ch:appx-lockstep}); fiducial-sensitivity
    figures and tables (Appendix~\ref{ch:appx-fiducials});
    PDF-reference manifest
    (Appendix~\ref{ch:appx-refs-manifest}); open-obligations
    roadmap (Appendix~\ref{ch:appx-roadmap}).
\end{description}

A reader interested only in the mass-gap statement and its proof
sketch can read Chapters~\ref{ch:introduction},
\ref{ch:sheaf-framework}, \ref{ch:sigma-hat-trajectory},
\ref{ch:continuum-measure}, and~\ref{ch:mass-gap-theorem} for a
$\sim 40$-page reduction.  All other chapters elaborate on this
core.
